\chapter{复变函数的积分}

\section{复变函数积分的定义}

\noindent \textbf{I. 复变函数积分的定义}

设 $C$ 是复平面上一条光滑（或逐段光滑）的曲线，我们选定一个正方向（规定 $+$：起点指向终点或逆时针），称 $C$ 为一条**有向曲线**，其反向曲线为 $C^-$。

\begin{definition}
设函数 $w=f(z)$ 在区域 $D$ 内有定义，$C$ 为 $D$ 内以 $A$ 为起点、$B$ 为终点的光滑有向曲线。
若将曲线 $C$ 任意划分为 $n$ 个小弧段，分点为 $A=z_0, z_1, \dots, z_n=B$。在每个小弧段 $\overset{\frown}{z_{k-1}z_k}$ 上任取一点 $\xi_k$，作和式：
\[
\sum_{k=1}^n f(\xi_k) \Delta z_k, \quad \text{其中 } \Delta z_k = z_k - z_{k-1}.
\]
记 $\delta = \max_{1\leq k\leq n} |\Delta z_k|$（即小弧段的最大长度）。若当 $\delta \to 0$ 时，该和式的极限存在且唯一（与分法、取点 $\xi_k$ 无关），则称此极限为函数 $f(z)$ 在曲线 $C$ 上的**复变函数积分**，记作
\[
\int_C f(z) \, dz = \lim_{\delta \to 0} \sum_{k=1}^n f(\xi_k) \Delta z_k.
\]
\end{definition}

\subsection{积分的计算方法（参数方程法）}
设曲线 $C$ 的参数方程为 $z = z(t) = x(t) + iy(t)$，其中 $t \in [\alpha, \beta]$。
则 $dz = z'(t) dt = (x'(t) + iy'(t)) dt$。
复积分可转化为定积分计算：
\[
\int_C f(z) dz = \int_\alpha^\beta f[z(t)] z'(t) dt.
\]
也可拆分为实部和虚部：
\[
\int_C f(z) dz = \int_C u dx - v dy + i \int_C v dx + u dy.
\]

\section{典型例题计算}

\begin{example}
计算 $\displaystyle \int_C f(z) dz$，其中 $f(z) = t + t^2 i$，路径 $C$ 分别为：
(1) 从点 $O$ 沿曲线 $z=t+it$ 到点 $1+i$；
(2) 从点 $O$ 沿实轴到 $1$，再平行虚轴到 $1+i$；
(3) $z \to 1$。

\textbf{解 (1)：} 路径 $z = t + it, \quad t: 0 \to 1$.
$dz = (1+i)dt$.
\[
\int_C f(z) dz = \int_0^1 (t + t^2 i)(1+i) dt = \int_0^1 (t - t^2) + i(3t^2) dt.
\]
\[
= \left( \frac{1}{2}t^2 - \frac{1}{3}t^3 \right)\bigg|_0^1 + i t^3 \bigg|_0^1 = \frac{1}{6} + i.
\]

\textbf{解 (2)：} 分为两段 $C_1: z=t, t:0\to 1$；$C_2: z=1+it, t:0\to 1$.
实部 $u=t, v=0$；虚部 $u=0, v=1$。
\[
\int_{C_1} (t+i^2)(1+2t+i) dt = \int_0^1 t dt = \frac{1}{2}.
\]
\[
\int_{C_2} (1+it+i^2)(i) dt = \int_0^1 (-t+i)i dt = \int_0^1 (-1 - ti) dt = -1 + \frac{1}{2}i.
\]
\[
\int_C z dz = \frac{1}{2} + (-1 + \frac{1}{2}i) = -\frac{1}{2} + i.
\]

\textbf{解 (3)：} $z \to 1$。
\[
\int_C z dz = \int_0^1 t dt = \frac{1}{2}.
\]
\end{example}

\begin{example}
设 $C$ 为正向圆周 $|z|=1$，计算：
(1) $\displaystyle \oint_C (x^2 - y^2 + 2xyi) dz$
(2) $\displaystyle \oint_C \frac{1}{z} dz$

\textbf{解：} 参数化 $z(\theta) = \cos\theta + i\sin\theta$，$\theta: 0 \to 2\pi$。
(1) $f(z) = (x^2-y^2) + i(2xy) = e^{i2\theta}$。
$dz = i e^{i\theta} d\theta$。
\[
\int_0^{2\pi} e^{i2\theta} \cdot i e^{i\theta} d\theta = i \int_0^{2\pi} e^{i3\theta} d\theta = 0.
\]
(2) $\displaystyle \oint_{|z|=1} \frac{1}{z} dz = \int_0^{2\pi} \frac{1}{e^{i\theta}} \cdot i e^{i\theta} d\theta = \int_0^{2\pi} i d\theta = 2\pi i$.
\end{example}

\section{重要公式与推论}

\subsection{复积分的估值不等式}
对于实变量积分，我们有：
\[
\left| \int_C f(z) dz \right| \le \int_C |f(z)| ds.
\]
若 $|f(z)| \le M$ 在 $C$ 上成立，$L$ 为 $C$ 的长度，则：
\[
\left| \int_C f(z) dz \right| \le M L.
\]

\subsection{计算 $\displaystyle \oint_{|z-z_0|=r} \frac{1}{(z-z_0)^n} dz$}
设 $C$ 为正向圆周 $|z-z_0|=r$ ($n$ 为正整数)。
令 $z-z_0 = re^{i\varphi}$，则 $dz = ire^{i\varphi} d\varphi$。
\[
f(z) = \frac{1}{(z-z_0)^n} = \frac{1}{r^n e^{in\varphi}} = \frac{1}{r^n} e^{-in\varphi}.
\]
\[
\oint_C \frac{1}{(z-z_0)^n} dz = \int_0^{2\pi} \frac{1}{r^n} e^{-in\varphi} \cdot ire^{i\varphi} d\varphi = \frac{i}{r^{n-1}} \int_0^{2\pi} e^{-i(n-1)\varphi} d\varphi.
\]
当 $n=1$ 时：
\[
\int_0^{2\pi} 1 \cdot d\varphi = 2\pi.
\]
当 $n \neq 1$ 时：
\[
\int_0^{2\pi} e^{-i(n-1)\varphi} d\varphi = 0.
\]
\textbf{结论：}
\[
\oint_{|z-z_0|=r} \frac{1}{(z-z_0)^n} dz =
\begin{cases}
2\pi i, & n=1 \\
0, & n \neq 1
\end{cases}
\]

\section{原函数与不定积分}

\begin{theorem}
设函数 $f(z)$ 在单连通区域 $B$ 内解析，则积分 $\displaystyle \int_C f(z) dz$ 与在 $B$ 内的路径无关，且 $F'(z) = f(z)$。
\end{theorem}

\begin{definition}
若 $F'(z) = f(z)$，则称 $F(z)$ 为 $f(z)$ 在区域 $B$ 内的一个**原函数**。
\end{definition}

\subsection{变上限积分的导数}
设 $F(z) = \int_{z_0}^z f(\zeta) d\zeta$，其中 $z_0$ 为 $B$ 内任意定点。
由解析性及积分与路径无关，可证：
\[
\lim_{\Delta z \to 0} \frac{F(z+\Delta z) - F(z)}{\Delta z} = f(z).
\]
即 $F'(z) = f(z)$。

\subsection{牛顿-莱布尼茨公式}
若 $G(z)$ 是 $f(z)$ 的一个原函数，则：
\[
\int_{z_0}^{z_1} f(z) dz = G(z_1) - G(z_0).
\]

\section{柯西-古萨基本定理}

\begin{theorem}[柯西-古萨定理]
如果函数 $w=f(z)$ 在单连通区域 $D$ 内解析，$C$ 是 $D$ 内的一条简单闭曲线，则
\[
\oint_C f(z) dz = 0.
\]
\end{theorem}

\subsection{闭路变形原理}
设 $f(z)$ 在区域 $D$ 内解析，$C, C_1, C_2, \dots, C_n$ 是 $D$ 内互不相交的正向简单闭曲线，且 $C_1, \dots, C_n$ 全部包含在 $C$ 所围的区域内，则
\[
\oint_C f(z) dz = \sum_{k=1}^n \oint_{C_k} f(z) dz.
\]

\section{复合闭路定理与例题}

\begin{example}
计算 $\displaystyle \oint_C \frac{3z-2i}{z(z-i)} dz$，其中 $C$ 为正向圆周 $|z|=2$。

\textbf{解：} 被积函数有两个奇点 $z=0$ 及 $z=i$，均在 $C$ 内部。
分式分解：
\[
\frac{3z-2i}{z(z-i)} = \frac{2}{z} + \frac{1}{z-i}.
\]
由闭路变形原理及柯西积分公式：
\[
\oint_C \frac{2}{z} dz = 2 \cdot 2\pi i = 4\pi i, \quad \oint_C \frac{1}{z-i} dz = 2\pi i.
\]
\[
\text{原式} = 4\pi i + 2\pi i = 6\pi i.
\]
\end{example}